\section{Avenues}

When the circuits get large and lengthy to run:

\subsection{Avenue 1: Machine Learning}

Trying to learn the relationship between the hidden bit strings (or possible ones) and the fucked up circuit in order to learn how the fucking up circuit is actually working. Purpose? When the circuits gets large we could maybe work out how to get the hidden bit strings (or most likely one) from the fucked up circuit.

\subsection{Avenue 2: Tensor Networks}

The MPS method forms a 1D chain. However, if the obfuscation creates nonlocal pairings, a tree can put paired qubits close and this could lead to better outcomes. So using the same MPO things.

How do we know this fails though

https://quantum-journal.org/papers/q-2023-03-30-964/

\subsection{Avenue 3: Circuit Compression}

Using a quantum computer
Using tensor networks - what tensor networks
Using the circuit compression

We need to try and reverse engineer to understand the obfuscation engine - look at the circuits and see if we can understand the obstruction engine. 

Learning some machine learning model to learn how the engine works based on the previous answers. For each challenge we will have multiple possible inputs.

What scales the best. How do the challenges progress and what works for one size may not work for a better.

As challenges progress qubit number increases.

Try different methods to find the best simulation method and compare on time and also amount of compression. If we want to use near term quantum computers the depth of the circuit matters a lot, but for classical simulation like MPS, the depth of the circuit isn't so important but the time it takes to run.


\subsection{Avenue 3: Make our own optimizer}
Maybe not because they give us a circuit compressor which is better.


\subsection{Avenue 4: Combine circuit compression and matrix product states}

Maybe compression to reduce entanglement gates and then MPS. The compression before could mean that the bond dimension doesn't get so high during the matrix product state approach and we don't have to make a cut.


\subsection{Avenue 5: Look at what QMill tools we have available to us - circuit compression, comparison etc.}

\subsection{Avenue 6: Look at the different methods in the paper as different ones are better for different types o circuits and use ML to pick}

A decision tree based on what kind of structure the circuit has.

\subsection{Avenue 7: Look at the 2D tensor network paper}

